\chapter{Semantics}
\label{ch:semantics}

This chapter gives DSTM a meaning in three layers.

\Cref{sec:l0} is the \emph{axiomatic} layer $\Lzm$: eight restrictions a
program must discharge, and one equation defining the store in terms of its
\gls{history}. It mentions no message, clock, version or distance. This is
what ``correct'' means.

\Cref{sec:l1} is the \emph{operational} layer $\Lom$: a labelled transition
system whose rules are the mechanism an implementation runs. Clocks, known
times, statuses, messages and termination appear here and nowhere else. A
refinement theorem justifies it against $\Lzm$.

\Cref{sec:sync} reads the same system as a dependence graph. Versioning is
renaming, distances are edge weights, drift is slack; acyclicity gives
deadlock freedom, the critical path gives the makespan, and the neighbourhood
diameter gives the memory bound.

\begin{table}[htbp]
\centering
\small
\begin{tabularx}{\textwidth}{@{}l l >{\raggedright\arraybackslash}X@{}}
\toprule
\textbf{Layer} & \textbf{Kind} & \textbf{Answers} \\
\midrule
$\Lzm$ (\cref{sec:l0})   & axiomatic   & What must a read return? \\
$\Lom$ (\cref{sec:l1})   & operational & By what mechanism, and when may it return? \\
$\Lim$ (implementation)  & code        & Over MPI, with what message types and buffers? \\
\bottomrule
\end{tabularx}
\caption{The three layers. This document specifies $\Lzm$ and $\Lom$ and
proves $\Lom \refines \Lzm$; the mapping from $\Lim$ to $\Lom$ is given by the
rule annotations in \cref{sec:l1-rules}, which cite the implementing source
line for each transition.}
\label{tab:layers}
\end{table}

\section{\texorpdfstring{$\Lzm$}{L0}: Axiomatic Semantics}
\label{sec:l0}

% ===========================================================================
\subsection{Histories}
\label{sec:l0-histories}

Everything here is a function of the \gls{history} (\cref{def:history}): the
\glspl{operation} issued, ordered per rank. No timing, message order or
interleaving is an input.

\begin{notationx}\label{not:history}
Throughout, $\Hist$ is a history over a rank set $\Ranks$ with a distance
\gls{comm-graph} $\cgraph$ (\cref{def:comm-graph}). We write $\Hist_{\opput}$ and
$\Hist_{\opget}$ for the \glspl{put} and \glspl{get} in $\Hist$, and
$\Hist|_p \defeq \{o \in \Hist : \opsrc{o} = p\}$ for the operations issued by
process $\proc{p}$, totally ordered by $\po$.
\end{notationx}

% ===========================================================================
\subsection{Restrictions}
\label{sec:l0-restrictions}

A history is not automatically meaningful. DSTM buys its simple memory model
by pushing the difficulty onto the program, as eight restrictions in three
groups:



\begin{description}[leftmargin=1.1cm,style=nextline,font=\normalfont\bfseries]
  \item[Semantic] \cref{res:write-disjoint} alone makes the store well
    defined. Without it ``the'' value of a \gls{cell} does not exist, and
    there is nothing to be correct \emph{about}.
  \item[Mechanical] \cref{res:monotone,res:lookahead,res:staleness,res:phase}
    make the store \emph{computable} by the mechanism of \cref{sec:l1}. A
    program violating them has a well-defined meaning DSTM cannot deliver.
  \item[Resource and lifecycle] \cref{res:horizon,res:collective,%
    res:terminate} bound memory and govern \gls{session} structure.
\end{description}

None of the eight is currently checked by the implementation; see
\crefrange{dev:zero-distance}{dev:barrier-doc}. \Cref{ch:restrictions}
collects them with proposed runtime checks.

\subsubsection*{Mechanical restrictions}

\begin{restr}[Monotone origin time]\label{res:monotone}
For every process $\proc{p}$ and every $o_1, o_2 \in \Hist|_p$ with $o_1 \po o_2$,
\[
  \optime{o_1} \le \optime{o_2}.
\]
Equivalently: a process never moves backwards along the timeline.
\devmark{dev:monotone-origin}
\end{restr}

\begin{restr}[Lookahead]\label{res:lookahead}
For every put $o \in \Hist_{\opput}$ with $p = \opsrc{o}$ and
$q = \optgt{o}$,
\[
  \opdelta{o} \;\ge\; \lag{p}{q},
  \qquad\text{and}\qquad
  q \in \nbr{p} \ \text{ or }\ q = p .
\]
The local case degenerates to $\opdelta{o} \ge 0$.
\devmark{dev:lookahead-uncheck}
\end{restr}

\begin{restr}[Staleness bound]\label{res:staleness}
For every get $o \in \Hist_{\opget}$ with $p = \opsrc{o}$ and
$q = \optgt{o}$,
\[
  -\opdelta{o} \;\le\; \lag{p}{q},
  \qquad\text{and}\qquad
  q \in \nbr{p} \ \text{ or }\ q = p .
\]
The local case degenerates to $\opdelta{o} \le 0$ combined with
$-\opdelta{o} \le 0$, hence $\opdelta{o} = 0$: a process reading its own \gls{region}
may only read its current effective time.
\devmark{dev:lookahead-uncheck}
\end{restr}

\Cref{res:lookahead,res:staleness} are duals, and reading them together is the
key to the whole model:

\begin{center}
\emph{The \gls{lag} is a lower bound on how long a write takes to
become visible,\\ and an upper bound on how far back a read may look.}
\end{center}

A write may take longer than $\lag{p}{q}$ and a read may look back less far;
what neither may do is beat the lag. Everything in \cref{sec:l1} is a
consequence.

\begin{restr}[Phase order]\label{res:phase}
For every process $\proc{p}$ and every $w \in \Hist_{\opput} \cap \Hist|_p$,
$g \in \Hist_{\opget} \cap \Hist|_p$ with
$\optgt{w} = \optgt{g}$ and $\opeff{w} = \opeff{g}$,
\[
  w \po g .
\]
That is: within one process, every write to a given version is issued before any
read of that version. \devmark{dev:phase-strictness}
\end{restr}

\Cref{res:phase} is stated per rank, which looks too weak: surely a
\emph{remote} write arriving after a read is served is equally fatal? It is,
but cannot happen.

\begin{proposition}[Minimality of the phase restriction]
\label{prop:phase-minimality}
Let $\Hist$ satisfy \cref{res:monotone,res:lookahead,res:phase}, and let
$g$ be a get targeting $(q,\eff)$ and $w$ a put targeting $(q,\eff)$ with
$\opsrc{w} \ne \opsrc{g}$. If $\opsrc{w} \ne q$ then
$\optime{w} \le \eff - \lag{\opsrc{w}}{q} < \eff$, whereas any operation
issued by $\opsrc{w}$ after a read of $(q,\eff)$ could be served has origin
time at least $\eff$. Hence the cross-process case of the phase order is implied,
and the per-rank statement of \cref{res:phase} is the only genuine
restriction.
\end{proposition}

\begin{proof}
By \cref{res:lookahead}, $\opdelta{w} \ge \lag{\opsrc{w}}{q} \ge 1$ using
(G3) and $\opsrc{w} \ne q$. Since $\opeff{w} = \optime{w} + \opdelta{w} =
\eff$, we get $\optime{w} \le \eff - 1 < \eff$. By \cref{res:monotone} the
\glspl{origin-time} of $\opsrc{w}$ are nondecreasing, so every put by $\opsrc{w}$
targeting $(q,\eff)$ occurs while that process's clock is still at most
$\eff - 1$. The publish guard of \cref{sec:l1-rules} does not fire until every
neighbour's \gls{known-time} $\kn{q}{r}$ satisfies $\kn{q}{r} + \lag{q}{r} > \eff$,
which for $r = \opsrc{w}$ requires $\opsrc{w}$ to have advanced strictly past
$\eff - \lag{q}{\opsrc{w}}$, i.e.\ past the last origin time at which it
could have produced such a put.

The remaining case is $\opsrc{w} = q \ne \opsrc{g}$, a local put on the target
process concurrent with a remote read. By \cref{res:staleness} applied to $g$,
$\eff \le \optime{g} - \lag{\opsrc{g}}{q} < \optime{g}$, and the read cannot
be served before $q$'s own clock reaches $\eff$; a local put at $\eff$ issued
after that point would violate \cref{res:monotone} unless it has
$\opdelta{} = 0$ and origin exactly $\eff$, which is the case
\cref{res:phase} governs on process $\proc{q}$ itself.
\end{proof}

This explains a puzzling feature of the implementation: the publish guard
tests $\clk{p} \ge \eff$, not $\clk{p} > \eff$
\devmark{dev:phase-strictness}. The non-strict test leaves exactly one hole
--- a local, zero-\gls{delta} write at the process's current time --- and
\cref{res:phase} plugs that hole and nothing more. A strict guard would close
it automatically but forbid a process from reading its own current tick, which
is the entire \glslink{push-model}{push idiom} (\cref{sec:exec-idioms}).

\subsubsection*{Semantic restriction}

\begin{restr}[Write disjointness]\label{res:write-disjoint}
For all distinct $w_1, w_2 \in \Hist_{\opput}$,
\[
  \fpw{w_1} \cap \fpw{w_2} = \emptyset .
\]
Equivalently: every cell is written at most once in the entire execution.
\devmark{dev:write-disjoint}
\end{restr}

Because a \gls{cell} carries an \gls{effective-time}, this is weaker than it
sounds: two writes to one address at different ticks do not conflict, only two
at the same tick. The \glslink{mailbox}{mailbox idiom}
(\cref{sec:gloss-idioms}) discharges it structurally.

\subsubsection*{Resource and lifecycle restrictions}

\begin{restr}[Bounded horizon]\label{res:horizon}
There is a constant $\delta_{\max}$ such that $\opdelta{o} \le \delta_{\max}$
for every $o \in \Hist_{\opput}$.
\end{restr}

This affects no result about values; only \cref{thm:space-bound} needs it. A
put with unbounded \gls{delta} creates a version that stays writable
indefinitely, and memory grows without limit.

\begin{restr}[Collective agreement]\label{res:collective}
Every collective entry point --- session start and end, the \gls{termination}
\gls{barrier}, and the creation of any derived structure --- is invoked by every
process, the same number of times, in the same order, with matching arguments.
\devmark{dev:noncollective}
\end{restr}

\begin{restr}[Termination bound]\label{res:terminate}
The argument $\ell$ passed to the terminating barrier satisfies
\[
  \ell \;\ge\; \max_{o \in \Hist} \opeff{o},
\]
and every process passes the same $\ell$. \devmark{dev:clock-conflation}
\end{restr}

\begin{definition}[Legality]\label{def:legal}
$\Legal{\Hist}$ holds when $\Hist$ satisfies
\crefrange{res:monotone}{res:terminate}.
\end{definition}

% ===========================================================================
\subsection{The store}
\label{sec:l0-store}

With legality in hand the semantics is a single equation.

\begin{definition}[Denotation of a history]\label{def:store-h}
For a history $\Hist$ satisfying \cref{res:write-disjoint}, the store
$\StoreH : \Cells \to \Bytes$ is
\[
  \StoreH\bigl(\cell{q}{\eff}{a}\bigr) \;\defeq\;
  \begin{cases}
    v_j
      & \text{if } \exists\, w = \Put{p}{q}{t}{\delta}{a_0}{v} \in \Hist
        \text{ with } t + \delta = \eff \\[-1pt]
      & \text{and } j = a - a_0 \in \{0, \dots, |v| - 1\}, \\[4pt]
    0 & \text{otherwise.}
  \end{cases}
\]
\end{definition}

\begin{lemma}[Well-definedness]\label{lem:store-welldef}
If $\Hist$ satisfies \cref{res:write-disjoint} then $\StoreH$ is a
well-defined total function.
\end{lemma}

\begin{proof}
Totality is immediate from the second clause. For single-valuedness, suppose
two puts $w_1 \ne w_2$ both satisfy the first clause for the cell
$c = \cell{q}{\eff}{a}$. Then $c \in \fpw{w_1}$ and $c \in \fpw{w_2}$ by
\cref{def:footprint}, contradicting \cref{res:write-disjoint}.
\end{proof}

\Cref{def:store-h} is the whole of the DSTM memory semantics: no clocks, no
lags, no \glspl{version}, no statuses, no messages. Those exist in
\cref{sec:l1} solely to compute this function without a global \gls{barrier}.

\begin{definition}[Read result]\label{def:read-result}
A get $g = \Get{p}{q}{t}{\delta}{a}{\ell}$ \emph{returns}
\[
  \StoreH(\cell{q}{t+\delta}{a})
  \;\cdots\;
  \StoreH(\cell{q}{t+\delta}{a+\ell-1}).
\]
\end{definition}

\begin{definition}[Default-valued read]\label{def:get-or-default}
A default-valued get carries an extra argument $d \in \Values$ with $|d| =
\ell$ and returns $d$ when $t + \delta < 0$, and behaves as
\cref{def:read-result} otherwise. This is the only construct that consults the
sign of an effective time; it exists so that the first ticks of a simulation,
whose producers have no history yet, need not be special-cased by the
application.
\end{definition}

\begin{definition}[Reads-from]\label{def:reads-from}
$w \rf g$ holds when $w \in \Hist_{\opput}$, $g \in \Hist_{\opget}$ and
$\fpw{w} \cap \fpr{g} \ne \emptyset$. By \cref{res:write-disjoint} each byte
returned by $g$ is supplied by at most one such $w$; bytes with no such $w$
read as $0$.
\end{definition}

% ===========================================================================
\subsection{Determinacy}
\label{sec:l0-determinacy}

$\StoreH$ depends on $\Hist$ alone, so fixing the \gls{history} fixes every
read. The question is whether the history itself is fixed: control flow
depends on what is read, which depends on what others wrote. The restrictions
make that circle well founded.

\begin{definition}[Deterministic program]\label{def:det-program}
A program is \emph{deterministic} when the operation each process issues next,
and the values it writes, are functions of that process's inputs and of the
values it has read so far.
\end{definition}

\begin{theorem}[Determinacy]\label{thm:determinism}
Let $P$ be a deterministic program all of whose executions are legal
(\cref{def:legal}). Then $P$ has a unique history $\Hist(P)$ and a unique
store $\Sigma_{\Hist(P)}$. In particular, every read returns the same value in
every execution, independently of scheduling, message order, and the relative
progress of processes.
\end{theorem}

\begin{proof}
Order operations by the lexicographic process
$\rho(o) \defeq \bigl(\opeff{o},\, \iota(o)\bigr)$, where $\iota(o)$ is the
position of $o$ in the \gls{program-order} of $\opsrc{o}$. We show that the value
returned by a get $g$ depends only on operations of strictly smaller $\rho$,
so that $\Hist(P)$ is determined by induction.

Let $g$ be a get targeting $(q,\eff)$ and let $w$ be any put with
$\fpw{w} \cap \fpr{g} \ne \emptyset$; then $\optgt{w} = q$ and
$\opeff{w} = \eff$, so the first components agree. Two cases.

\emph{Case $\opsrc{w} \ne \opsrc{g}$.} By \cref{prop:phase-minimality}, $w$ is
issued by a process whose clock is strictly below $\eff$ at the time, while $g$'s
issuer has clock at least $\eff$; the two are ordered by the causal chain of
\glspl{heartbeat}, and $w$'s issue necessarily precedes $g$'s completion. Its value
therefore depends only on reads that $\opsrc{w}$ performed at origin times
$\le \eff - \lag{\opsrc{w}}{q} < \eff$, all of strictly smaller first
component by (G3).

\emph{Case $\opsrc{w} = \opsrc{g} = p$.} Then $\optgt{w} = q$ and
$\opeff{w} = \opeff{g}$, so \cref{res:phase} gives $w \po g$, hence
$\iota(w) < \iota(g)$ and $\rho(w) < \rho(g)$.

In both cases every put contributing to $g$'s result has strictly smaller
$\rho$. The relation is well founded because \glspl{effective-time} are bounded below
by $0$ on any version that exists and program-order positions are natural
numbers. Determinism of $P$ then fixes each successive operation, and
\cref{lem:store-welldef} fixes each read result.
\end{proof}

\begin{corollary}[Equivalence with lockstep execution]\label{thm:bsp}
Let $P$ be as in \cref{thm:determinism}, and let $\mathrm{BSP}(P)$ be the
execution in which every process completes tick $t$ and passes a global barrier
before any process begins tick $t+1$. Then $\mathrm{BSP}(P)$ is a legal execution
of $P$ and produces $\Sigma_{\Hist(P)}$. Consequently every DSTM execution of
$P$ returns exactly the values the barrier-synchronised execution would.
\end{corollary}

\begin{proof}
$\mathrm{BSP}(P)$ satisfies every restriction, since a global barrier at each
tick boundary only adds ordering. By \cref{thm:determinism} the history is
unique, so $\mathrm{BSP}(P)$ realises $\Hist(P)$.
\end{proof}

\Cref{thm:bsp} licenses the whole design: removing the per-tick \gls{barrier}
and letting \glspl{process} drift changes \emph{when} work happens and
\emph{how much memory} is live, provably not \emph{what is computed}. Unlike
optimistic simulation nothing is ever discarded, because nothing is ever
wrong.

% ===========================================================================
\subsection{Objects as transition systems}
\label{sec:l0-ts}

The \gls{store} is static; what moves are the clocks, known times and
\glspl{version-status}. Each is a transition system, and \cref{sec:l1}'s
global system is their composition. Separating them shows which invariants are
local and which belong to the composition.

\begin{definition}[Transition system]\label{def:ts}
A \emph{transition system} is a tuple
$\TS{} = \langle \States, \Init, \Labels, \steps \rangle$ with $\States$ a set
of states, $\Init \in \States$ the initial state, $\Labels$ a set of labels
and ${\steps} \subseteq \States \times \Labels \times \States$. We write
$s \stepl{\lambda} s'$ for $(s,\lambda,s') \in {\steps}$.
\end{definition}

\begin{definition}[Clock]\label{ts:clock}
$\TS{\mathrm{clk}}(p) \defeq \langle \Time,\; 0,\;
 \{\optick(t) : t \in \Time\},\; \steps \rangle$ with
\[
  c \stepl{\optick(t)} t \quad\text{iff}\quad t > c .
\]
Monotone and deterministic in the label.
\end{definition}

\begin{definition}[Known time]\label{ts:known}
$\TS{\mathrm{kn}}(p,q) \defeq \langle \Time,\; 0,\;
 \{\mathsf{hear}(t) : t \in \Time\},\; \steps \rangle$ with
\[
  k \stepl{\mathsf{hear}(t)} \max(k, t).
\]
Total --- every label is enabled in every state --- and monotone.
\end{definition}

\begin{definition}[Version status]\label{ts:version}
$\TS{\mathrm{st}}(p,\eff) \defeq \langle \Status,\; \Absent,\;
 \{\mathsf{create}, \mathsf{publish}, \mathsf{retire}, \mathsf{evict}\},\;
 \steps \rangle$ with
\[
  \Absent \stepl{\mathsf{create}} \Wr
  \stepl{\mathsf{publish}} \Rd
  \stepl{\mathsf{retire}} \Dc
  \stepl{\mathsf{evict}} \Absent .
\]
Every transition is monotone in the order of \cref{def:status} except
$\mathsf{evict}$, which returns the system to $\Absent$. That single
non-monotone edge is why \cref{thm:discard-safety} is needed: it must be
impossible to distinguish an evicted version from one that never existed,
which holds exactly when nothing will ever ask for it again.
\end{definition}

\begin{definition}[Cell]\label{ts:cell}
$\TS{\mathrm{cell}}(c) \defeq \langle \Bytes,\; 0,\;
 \{\mathsf{write}(v) : v \in \Bytes\},\; \steps \rangle$ with
$b \stepl{\mathsf{write}(v)} v$, subject to the guard that
$\mathsf{write}$ is enabled only while the enclosing version's status is
$\Wr$. Under \cref{res:write-disjoint} at most one $\mathsf{write}$ ever
occurs, so the reachable states are exactly $\{0\} \cup \{v\}$ and the
system is confluent.
\end{definition}

\begin{definition}[Allocation]\label{ts:alloc}
$\TS{\mathrm{alloc}}(p) \defeq \langle \States,\; \emptyset,\;
 \{\mathsf{alloc}(n,h,I), \mathsf{free}(h)\},\; \steps \rangle$, where
$\States$ is the set of well-formed allocation states
(\cref{def:allocation}) and
\[
  \allocof{p} \stepl{\mathsf{alloc}(n,h,I)} \allocof{p}[\,I \mapsto h\,],
  \qquad
  \allocof{p} \stepl{\mathsf{free}(h)} \restrict{\allocof{p}}{\Addrof{p}
    \setminus \extent{p}{h}},
\]
the first enabled when $h$ is fresh and $I$ is a length-$n$ interval disjoint
from $\live{p}$. The label carries $I$: the system is \emph{nondeterministic},
and every resolution of that nondeterminism is a conforming allocator
(\cref{rem:alloc-policy}).
\end{definition}

\begin{remark}[Where the difficulty lives]
Each of \crefrange{ts:clock}{ts:alloc} is trivial in isolation: two monotone
counters, a four-state chain, a single-assignment register and a disjointness
bookkeeper. All the difficulty is in the \emph{guards} coupling them --- when
$\mathsf{publish}$ and $\mathsf{retire}$ may fire, which is a question about
other processes. \Cref{sec:l1} supplies them.
\end{remark}

\Cref{ts:alloc} is the one object with no such guard: allocation is purely
local, involves no message and constrains no other process. That is why it
does not appear in \cref{sec:l1-rules} --- its work is done before any
operation is issued.

% ===========================================================================
\subsection{The correctness criterion}
\label{sec:l0-correctness}

\begin{definition}[Correct implementation]\label{def:correctness}
An implementation is \emph{correct} if, for every deterministic program $P$
all of whose executions are legal, every execution of $P$
\begin{enumerate}[label=(\roman*),leftmargin=1.4cm]
  \item terminates, with every get having completed, and
  \item returns, for each get, the bytes prescribed by
    \cref{def:read-result} with respect to $\Sigma_{\Hist(P)}$.
\end{enumerate}
\end{definition}

Clause (ii) is safety and is discharged by \cref{thm:read-soundness}; clause
(i) is liveness and is discharged by \cref{thm:deadlock-freedom} together with
\cref{thm:barrier-correct}. Both are proved against $\Lom$ in the next
section.

\section{\texorpdfstring{$\Lom$}{L1}: Operational Semantics}
\label{sec:l1}

$\Lzm$ said what a read must return, not how a process could know, without
asking anyone, that a version is complete. That is $\Lom$: a labelled
transition system whose rules are the mechanism the runtime executes.

Each rule is annotated with the source location implementing it; those
annotations are the refinement map from $\Lim$ to $\Lom$.
\Cref{sec:l1-refinement} discharges $\Lom$ against $\Lzm$.

% ===========================================================================
\subsection{Configurations}
\label{sec:l1-config}

\begin{definition}[Local state]\label{def:local-state}
The state of process $\proc{p}$ is a quadruple $(\clk{p}, K_p, \Vers{p}, \Pend{p})$:
\begin{itemize}[leftmargin=1.5em,itemsep=1pt]
  \item $\clk{p} \in \Time$, the \gls{local-clock} (\cref{ts:clock});
  \item $K_p : \Ranks \to \Time$, the \glspl{known-time} (\cref{ts:known}), written
    $\kn{p}{q}$;
  \item $\Vers{p} : \Time \pto \Status \times (\Addrof{p} \to \Bytes)$, the
    partial map of live \glspl{version}, with components $\Vst{p}{\eff}$ and
    $\Vbuf{p}{\eff}$; we set $\Vst{p}{\eff} = \Absent$ when
    $\eff \notin \operatorname{dom} \Vers{p}$;
  \item $\Pend{p}$, a finite multiset of pending requests.
\end{itemize}
\end{definition}

\begin{definition}[Requests]\label{def:requests}
A pending request is one of
\[
\begin{array}{ll}
  \mathsf{lput}(\eff, a, v)            & \text{local put, not yet applied}\\
  \mathsf{lget}(\eff, a, \ell)         & \text{local get, awaiting publication}\\
  \mathsf{iput}(\eff, a, v)            & \text{received put, not yet applied}\\
  \mathsf{iget}(\eff, a, \ell, r, \iota) & \text{received read request from } r\\
  \mathsf{oput}                        & \text{outgoing put, awaiting send completion}\\
  \mathsf{oget}(\eff, \iota)           & \text{outgoing read, awaiting reply}
\end{array}
\]
We write $\opeff{\varrho}$ for the \gls{effective-time} recorded in request
$\varrho$. The kinds mirror \srcrefr{request.h}{33}{40} exactly.
\end{definition}

\begin{definition}[Messages]\label{def:messages}
$\Net$ is a finite multiset of in-flight messages, each of one of the five
forms of \srcrefr{msg.h}{39}{43}:
\[
\begin{array}{ll}
  \mathsf{PUT}(p \to q,\, t,\, \delta,\, a,\, v) &
  \mathsf{GETREQ}(p \to q,\, t,\, \delta,\, a,\, \ell,\, \iota) \\
  \mathsf{GETREP}(q \to p,\, \eff,\, v,\, \iota) &
  \mathsf{HB}(p \to q,\, t) \\
  \mathsf{TOK}(p \to q,\, \ell) &
\end{array}
\]
Every message carries a \emph{stamp}: the value $t$ for $\mathsf{PUT}$,
$\mathsf{GETREQ}$ and $\mathsf{HB}$, the value $\eff$ for $\mathsf{GETREP}$,
and $\ell$ for $\mathsf{TOK}$. We write $\mathit{stamp}(m)$ for it.
\end{definition}

\begin{definition}[Global configuration]\label{def:configuration}
$\Cfg \defeq \bigl\langle (\clk{p}, K_p, \Vers{p}, \Pend{p})_{p \in \Ranks},\
\Net \bigr\rangle$. The initial configuration $\Cfg_0$ has $\clk{p} = 0$,
$\kn{p}{q} = 0$, $\Vers{p} = \emptyset$, $\Pend{p} = \emptyset$ and
$\Net = \emptyset$ for all $p, q$.
\end{definition}

\begin{convention}[Notation in rules]\label{conv:rule-notation}
$\Cfg[\,x \mapsto e\,]$ denotes $\Cfg$ with component $x$ replaced by $e$. We
write $\Pend{p} \uplus \{\varrho\}$ for multiset insertion and
$\Net \uplus \{m\}$ for message injection. Rules are read as: if the premises
hold in $\Cfg$, then $\Cfg \steps \Cfg'$ with $\Cfg'$ as described in the
conclusion.
\end{convention}

% ===========================================================================
\subsection{Transition rules}
\label{sec:l1-rules}

\subsubsection*{Clock and publication of progress}

\begin{mathpar}
\inferrule*[left=\rn{Tick}]
  {t > \clk{p}}
  {\clk{p} \mathrel{{:}{=}} t \\\\
   \Net \mathrel{{:}{=}} \Net \uplus
     \{\, \mathsf{HB}(p \to q,\, t) : q \in \nbr{p} \,\}}
\end{mathpar}

\noindent\srcrefr{ops.c}{37}{53} (broadcast), \srcref{ops.c}{603} and
\srcref{ops.c}{672} (the guard, folded into put and get issue). \rn{Tick}
is not a separate API call: it fires as a side effect of issuing an \gls{operation}
whose \gls{origin-time} exceeds the clock.

\subsubsection*{Issuing operations}

\begin{mathpar}
\inferrule*[left=\rn{Put-Local}]
  {\delta \ge 0 \\ q = p}
  {\Pend{p} \mathrel{{:}{=}} \Pend{p} \uplus
     \{\mathsf{lput}(t + \delta,\, a,\, v)\}}

\inferrule*[left=\rn{Put-Remote}]
  {\delta \ge 1 \\ q \in \nbr{p}}
  {\Pend{p} \mathrel{{:}{=}} \Pend{p} \uplus \{\mathsf{oput}\} \\\\
   \Net \mathrel{{:}{=}} \Net \uplus
     \{\mathsf{PUT}(p \to q,\, t,\, \delta,\, a,\, v)\}}
\end{mathpar}

\begin{mathpar}
\inferrule*[left=\rn{Get-Local}]
  {\delta = 0 \\ q = p}
  {\Pend{p} \mathrel{{:}{=}} \Pend{p} \uplus
     \{\mathsf{lget}(t,\, a,\, \ell)\}}

\inferrule*[left=\rn{Get-Remote}]
  {\delta \le -1 \\ q \in \nbr{p} \\ \iota \text{ fresh}}
  {\Pend{p} \mathrel{{:}{=}} \Pend{p} \uplus \{\mathsf{oget}(t+\delta, \iota)\} \\\\
   \Net \mathrel{{:}{=}} \Net \uplus
     \{\mathsf{GETREQ}(p \to q,\, t,\, \delta,\, a,\, \ell,\, \iota)\}}
\end{mathpar}

\noindent\srcrefr{ops.c}{608}{636} and \srcrefr{ops.c}{677}{720}. The
$\delta = 0 \Rightarrow q = p$ side conditions are the \code{-DSTM\_EDOM}
checks at \srcref{ops.c}{598} and \srcref{ops.c}{667}. Note what the rules do
\emph{not} require: neither checks $\delta$ against $\dist{p}{q}$. That is
\cref{res:lookahead,res:staleness}, and it is assumed, not verified
\devmark{dev:lookahead-uncheck}.

\subsubsection*{Message delivery}

Every delivery raises the receiver's known time; this is the only way known
times ever change.

\begin{mathpar}
\inferrule*[left=\rn{Recv}]
  {m \in \Net \\ \mathit{dst}(m) = p \\ \mathit{src}(m) = r}
  {\Net \mathrel{{:}{=}} \Net \setminus \{m\} \\\\
   \kn{p}{r} \mathrel{{:}{=}} \max\bigl(\kn{p}{r},\, \mathit{stamp}(m)\bigr) \\\\
   \rn{Dispatch}(m)}
\end{mathpar}

\noindent\srcrefr{ops.c}{432}{436}. $\rn{Dispatch}$ is:

\begin{mathpar}
\inferrule*[left=\rn{D-Put}]
  {m = \mathsf{PUT}(r \to p,\, t,\, \delta,\, a,\, v)}
  {\Pend{p} \mathrel{{:}{=}} \Pend{p} \uplus \{\mathsf{iput}(t+\delta,\, a,\, v)\}}

\inferrule*[left=\rn{D-Req}]
  {m = \mathsf{GETREQ}(r \to p,\, t,\, \delta,\, a,\, \ell,\, \iota)}
  {\Pend{p} \mathrel{{:}{=}} \Pend{p} \uplus
     \{\mathsf{iget}(t+\delta,\, a,\, \ell,\, r,\, \iota)\}}
\end{mathpar}

\begin{mathpar}
\inferrule*[left=\rn{D-Rep}]
  {m = \mathsf{GETREP}(r \to p,\, \eff,\, v,\, \iota) \\
   \mathsf{oget}(\eff, \iota) \in \Pend{p}}
  {\Pend{p} \mathrel{{:}{=}} \Pend{p} \setminus \{\mathsf{oget}(\eff,\iota)\} \\\\
   \text{the issuing get returns } v}

\inferrule*[left=\rn{D-Beat}]
  {m = \mathsf{HB}(r \to p,\, t)}
  {\text{no further effect}}
\end{mathpar}

\noindent\srcrefr{ops.c}{438}{506}. \rn{D-Beat} has no effect beyond the
known-time update performed by \rn{Recv}: that update is the heartbeat's
entire purpose.

\subsubsection*{Applying writes}

\begin{mathpar}
\inferrule*[left=\rn{Apply}]
  {\varrho \in \Pend{p} \\
   \varrho \in \{\mathsf{lput}(\eff,a,v),\, \mathsf{iput}(\eff,a,v)\} \\
   \Vst{p}{\eff} \in \{\Absent, \Wr\}}
  {\Vst{p}{\eff} \mathrel{{:}{=}} \Wr \\\\
   \Vbuf{p}{\eff} \mathrel{{:}{=}} \Vbuf{p}{\eff}[a \dots a{+}|v|{-}1 \mapsto v] \\\\
   \Pend{p} \mathrel{{:}{=}} \Pend{p} \setminus \{\varrho\}}
\end{mathpar}

\noindent\srcrefr{ops.c}{223}{258} (local) and \srcrefr{ops.c}{282}{315}
(incoming). Creating a version initialises $\Vbuf{p}{\eff}$ to all zeroes
(\cref{conv:no-inherit}), implemented as the \code{calloc} in
\code{dstm\_region\_version\_find\_or\_create}.

\begin{mathpar}
\inferrule*[left=\rn{Apply-Violation}]
  {\varrho \in \Pend{p} \\
   \varrho \in \{\mathsf{lput}(\eff,a,v),\, \mathsf{iput}(\eff,a,v)\} \\
   \Vst{p}{\eff} \in \{\Rd, \Dc\}}
  {\textbf{abort}}
\end{mathpar}

\noindent\srcrefr{ops.c}{244}{253} and \srcrefr{ops.c}{301}{310}.
\rn{Apply-Violation} is not a semantic rule: it is a detector. On a legal
history it is unreachable (\cref{thm:write-closure}), so reaching it proves
that \cref{res:phase}, \cref{res:monotone} or \cref{res:lookahead} was
violated. The implementation deliberately aborts rather than dropping or
applying the write, because both alternatives silently corrupt a value some
reader has already observed.

\subsubsection*{Promotion --- the two rules that matter}

\begin{mathpar}
\inferrule*[left=\rn{Publish}]
  {\Vst{p}{\eff} = \Wr \\
   \nexists\, \varrho \in \Pend{p}.\;
     \varrho \in \{\mathsf{lput},\mathsf{iput}\} \wedge \opeff{\varrho} = \eff \\\\
   \forall q \in \nbr{p}.\;\ \kn{p}{q} + \dist{p}{q} > \eff \\
   \clk{p} \ge \eff}
  {\Vst{p}{\eff} \mathrel{{:}{=}} \Rd}
\end{mathpar}

\noindent\srcrefr{ops.c}{57}{117}: the pending-put scan at
\srcrefr{ops.c}{63}{76}, the neighbour scan at \srcrefr{ops.c}{86}{101}, the
local-clock clause at \srcref{ops.c}{112}.

\begin{mathpar}
\inferrule*[left=\rn{Retire}]
  {\Vst{p}{\eff} = \Rd \\
   \nexists\, \varrho \in \Pend{p}.\;
     \varrho \in \{\mathsf{lget},\mathsf{iget}\} \wedge \opeff{\varrho} = \eff \\\\
   \forall q \in \nbr{p}.\;\ \kn{p}{q} - \dist{p}{q} > \eff \\
   \clk{p} \ge \eff}
  {\Vst{p}{\eff} \mathrel{{:}{=}} \Dc}
\end{mathpar}

\noindent\srcrefr{ops.c}{121}{172}.

\begin{mathpar}
\inferrule*[left=\rn{Evict}]
  {\Vst{p}{\eff} = \Dc}
  {\Vers{p} \mathrel{{:}{=}} \Vers{p} \setminus \{\eff\}}
\end{mathpar}

\noindent\srcrefr{ops.c}{201}{205}.

The two neighbour clauses are the design in one line. A neighbour $q$ believed
to be at $\kn{p}{q}$ can still write as early as $\kn{p}{q} + \dist{p}{q}$ and
still read as far back as $\kn{p}{q} - \dist{p}{q}$. Publication waits for the
first bound to pass $\eff$, retirement for the second: one \gls{distance}
serving as a forward horizon for writes and a backward horizon for reads.

\subsubsection*{Serving reads}

\begin{mathpar}
\inferrule*[left=\rn{Serve-Local}]
  {\mathsf{lget}(\eff, a, \ell) \in \Pend{p} \\ \Vst{p}{\eff} = \Rd}
  {\Pend{p} \mathrel{{:}{=}} \Pend{p} \setminus \{\mathsf{lget}(\eff,a,\ell)\} \\\\
   \text{the issuing get returns } \Vbuf{p}{\eff}[a \dots a{+}\ell{-}1]}
\end{mathpar}

\begin{mathpar}
\inferrule*[left=\rn{Serve-Remote}]
  {\mathsf{iget}(\eff, a, \ell, r, \iota) \in \Pend{p} \\ \Vst{p}{\eff} = \Rd}
  {\Pend{p} \mathrel{{:}{=}} \Pend{p} \setminus \{\mathsf{iget}(\dots)\} \\\\
   \Net \mathrel{{:}{=}} \Net \uplus
     \{\mathsf{GETREP}(p \to r,\, \eff,\, \Vbuf{p}{\eff}[a \dots a{+}\ell{-}1],\, \iota)\}}
\end{mathpar}

\noindent\srcrefr{ops.c}{260}{280} and \srcrefr{ops.c}{317}{357}. Both rules
require status exactly $\Rd$: a $\Wr$ version is incomplete and a $\Dc$ one is
already forfeit. A get whose target version is $\Absent$ simply does not
fire, which is why violating \cref{res:staleness} manifests as a hang rather
than an error \devmark{dev:lookahead-uncheck}.

\subsubsection*{Termination}

The \gls{barrier} is a three-round token ring, summarised as three rules; the
five-state machine implementing them is \srcrefr{barrier.c}{113}{225}.

\begin{mathpar}
\inferrule*[left=\rn{Bar-One}]
  {p \text{ has no } \mathsf{oput}/\mathsf{oget} \text{ pending} \\
   p = 0 \vee \text{first-round token received}}
  {\text{forward first-round token to } p \oplus 1}

\inferrule*[left=\rn{Bar-Two}]
  {\text{first ring closed} \\ \Pend{p} = \emptyset}
  {\mathit{drained}_p \mathrel{{:}{=}} \mathsf{true} \\\\
   \text{forward second-round token to } p \oplus 1}
\end{mathpar}

\begin{mathpar}
\inferrule*[left=\rn{Bar-Exit}]
  {\text{second ring closed}}
  {\text{forward exit token} \\\\ \mathit{drained}_p \mathrel{{:}{=}} \mathsf{false} \\\\
   \text{barrier returns}}
\end{mathpar}

While $\mathit{drained}_p$ holds, \rn{Publish} and \rn{Retire} drop their
neighbour clauses (\srcref{ops.c}{85}, \srcref{ops.c}{149}). This is sound
because the closed first ring certifies no process will issue another
\gls{operation}, making those horizons vacuous; and necessary, since without
it the last \glspl{version} could never publish, no further \glspl{heartbeat}
arriving to raise the known times.

% ===========================================================================
\subsection{Fairness}
\label{sec:l1-fairness}

Safety holds for every execution; liveness needs the usual assumptions for an
asynchronous message-passing system.

\begin{definition}[Admissible execution]\label{def:admissible}
An execution is \emph{admissible} when
\begin{enumerate}[label=(F\arabic*),leftmargin=1.5cm,itemsep=1pt]
  \item \label{f:delivery} every message in $\Net$ is eventually delivered by
    \rn{Recv} (reliable, eventual, no loss or duplication);
  \item \label{f:progress} every rule other than \rn{Tick} that stays
    continuously enabled eventually fires --- in particular a process blocked in
    a get keeps driving the engine (\srcref{ops.c}{724});
  \item \label{f:clock} every process either eventually enters the terminating
    barrier or advances its clock without bound.
\end{enumerate}
\end{definition}

\ref{f:clock} formalises the warning in \cref{sec:exec-clocks} that a silent
process blocks its neighbours. It is a restriction on the \emph{application},
not the runtime.

% ===========================================================================
\subsection{Invariants}
\label{sec:l1-invariants}

\begin{invariant}[Clock soundness]\label{inv:l1-clock}
$\kn{p}{q} \le \clk{q}$ for all $p,q$, at every reachable configuration.
\end{invariant}

\begin{proof}
Initially both sides are $0$. $\kn{p}{q}$ changes only in \rn{Recv}, to
$\mathit{stamp}(m)$ for a message $m$ sent by $q$. Each message form is
stamped at or below $q$'s clock at send time: $\mathsf{PUT}$, $\mathsf{GETREQ}$
and $\mathsf{HB}$ carry the origin time $t$, and \rn{Tick} has already set
$\clk{q} \ge t$; $\mathsf{GETREP}$ carries $\eff = t + \delta \le t \le \clk{q}$
because \glspl{get} have $\delta \le 0$; $\mathsf{TOK}$ carries $\ell$, which by
\cref{res:terminate} bounds every effective time and is passed identically by
all \glspl{rank}. Since $\clk{q}$ is nondecreasing (\cref{ts:clock}), the bound
survives.
\end{proof}

\Cref{prop:conservative} now has teeth: known times are lower bounds, both
guards are monotone in them, so every promotion the runtime makes is one a
perfectly-informed oracle would also make. The runtime is never wrong, only
sometimes late.

% ===========================================================================
\subsection{Safety}
\label{sec:l1-safety}

\begin{theorem}[Write closure]\label{thm:write-closure}
Let $\Hist$ be legal. If $\Vst{p}{\eff} \ge \Rd$ in some reachable
configuration, then no operation issued at or after that point is a \gls{put}
targeting $(p, \eff)$, and $\Pend{p}$ contains no unapplied put for $\eff$.
Consequently \rn{Apply-Violation} is unreachable.
\end{theorem}

\begin{proof}
The second clause is the first premise of \rn{Publish}. For the first,
consider a put $w$ targeting $(p,\eff)$ issued by process $\proc{r}$ after the
promotion.

\emph{Case $r \ne p$.} Then $r \in \nbr{p}$, since \cref{res:lookahead}
requires $r \in \nbr{p}$. The \rn{Publish} premise gives
$\kn{p}{r} + \dist{p}{r} > \eff$, and \cref{inv:l1-clock} gives
$\clk{r} \ge \kn{p}{r}$, hence
\[
  \clk{r} + \dist{r}{p} \;\ge\; \kn{p}{r} + \dist{p}{r} \;>\; \eff,
\]
using (M3) for $\dist{r}{p} = \dist{p}{r}$. By \cref{res:monotone} the origin
time of $w$ is at least $\clk{r}$ as it stands now, and by
\cref{res:lookahead} $\opdelta{w} \ge \dist{r}{p}$. Therefore
$\opeff{w} = \optime{w} + \opdelta{w} \ge \clk{r} + \dist{r}{p} > \eff$, so
$w$ does not target $\eff$. Contradiction.

\emph{Case $r = p$.} The \rn{Publish} premise gives $\clk{p} \ge \eff$. By
\cref{res:monotone}, $\optime{w} \ge \clk{p} \ge \eff$, and
$\opdelta{w} \ge 0$, so $\opeff{w} \ge \eff$ with equality only if
$\optime{w} = \eff$ and $\opdelta{w} = 0$. That residual case is exactly a
local zero-\gls{delta} write at the process's current time, and it is excluded by
\cref{res:phase}: the promotion to $\Rd$ can only have been triggered by a get
on $(p,\eff)$ having been issued, and \cref{res:phase} forbids $p$ from
issuing a put on $(p,\eff)$ after such a get.
\end{proof}

\Cref{thm:write-closure} consumes all four mechanical restrictions, and is
where (M2) and (M3) earn their place. Dropping symmetry breaks the inequality
chain in the first case \devmark{dev:asymmetric}; dropping (M2) makes the
second case's strictness fail for a neighbour at distance zero
\devmark{dev:zero-distance}.

\begin{corollary}[Immutability]\label{thm:immutability}
For a legal history, $\Vbuf{p}{\eff}$ is constant from the moment
$\Vst{p}{\eff}$ reaches $\Rd$ until the version is evicted.
\end{corollary}

\begin{proof}
$\Vbuf{p}{\eff}$ is modified only by \rn{Apply}, whose premise requires
$\Vst{p}{\eff} \in \{\Absent, \Wr\}$. By \cref{thm:write-closure} no such
request exists after publication, and statuses never decrease except through
\rn{Evict}, which removes the version outright.
\end{proof}

\begin{theorem}[Discard safety]\label{thm:discard-safety}
Let $\Hist$ be legal. If $\Vst{p}{\eff} = \Dc$ in some reachable
configuration, then no operation issued at or after that point is a get
targeting $(p,\eff)$.
\end{theorem}

\begin{proof}
Let $g$ be such a get, issued by $r$.

\emph{Case $r \ne p$.} Then $r \in \nbr{p}$ by \cref{res:staleness}. The
\rn{Retire} premise gives $\kn{p}{r} - \dist{p}{r} > \eff$, and
\cref{inv:l1-clock} gives $\clk{r} \ge \kn{p}{r}$, so
$\clk{r} - \dist{r}{p} > \eff$ by (M3). By \cref{res:monotone},
$\optime{g} \ge \clk{r}$; by \cref{res:staleness},
$\opdelta{g} \ge -\dist{r}{p}$. Hence
$\opeff{g} \ge \clk{r} - \dist{r}{p} > \eff$.

\emph{Case $r = p$.} The premise gives $\clk{p} \ge \eff$, and
\cref{res:staleness} forces $\opdelta{g} = 0$ for a local get, so
$\opeff{g} = \optime{g} \ge \clk{p} \ge \eff$. Equality would require
$\optime{g} = \eff = \clk{p}$; but retirement additionally requires no pending
$\mathsf{lget}$ at $\eff$, and by \cref{res:phase} combined with
\cref{thm:write-closure} the version was already published, so any local get
at $\eff$ was issued before retirement and has been served.
\end{proof}

\begin{theorem}[Read soundness: $\Lom \refines \Lzm$]\label{thm:read-soundness}
Let $\Hist$ be a legal history. Every get in $\Hist$ that completes in an
$\Lom$ execution returns exactly the bytes prescribed by
\cref{def:read-result} with respect to $\StoreH$.
\end{theorem}

\begin{proof}
A get on $(q,\eff)$ returns via \rn{Serve-Local} or \rn{Serve-Remote},
both of which read $\Vbuf{q}{\eff}$ while $\Vst{q}{\eff} = \Rd$. Define the
abstraction $\alpha(\Cfg)$ mapping each published or retired version to the
byte function $\Vbuf{q}{\eff}$. We show $\alpha(\Cfg)$ agrees with $\StoreH$
on every published version.

Fix a \gls{cell} $c = \cell{q}{\eff}{a}$ with $\Vst{q}{\eff} = \Rd$. By
\cref{thm:write-closure}, every put in $\Hist$ targeting $c$ has already been
applied by \rn{Apply}. By \cref{res:write-disjoint} there is at most one
such put $w$, so no application overwrote another, and the byte in
$\Vbuf{q}{\eff}$ at offset $a$ is the byte $w$ supplied --- which is
$\StoreH(c)$ by \cref{def:store-h}. If there is no such put, no
\rn{Apply} ever touched offset $a$ of this version, so it retains the
zero from version creation (\cref{conv:no-inherit}), and
$\StoreH(c) = 0$ by the second clause of \cref{def:store-h}.
By \cref{thm:immutability} the value does not subsequently change, so every
get on $c$ returns the same byte.
\end{proof}

% ===========================================================================
\subsection{Liveness}
\label{sec:l1-liveness}

Safety says a completed get is right; liveness says it completes.

\begin{theorem}[Deadlock freedom]\label{thm:deadlock-freedom}
Let $\Hist$ be legal and let the execution be admissible
(\cref{def:admissible}). Then every get completes.
\end{theorem}

\begin{proof}
Suppose not, and consider a get blocked forever on $(q, \eff)$. It is waiting
for \rn{Publish} at $(q,\eff)$, whose premises are of three kinds.

The pending-put clause discharges itself: $\Pend{q}$ is finite, no new put
targeting $\eff$ can be created once $q$'s neighbours have advanced (see
below), and \rn{Apply} is always enabled for a pending put on a
non-published version, so by \ref{f:progress} they all clear.

The local-clock clause $\clk{q} \ge \eff$ discharges by \ref{f:clock}: $q$
either advances past $\eff$ or enters the barrier, and in the barrier
$\mathit{drained}_q$ removes the clause.

The neighbour clause requires, for each $r \in \nbr{q}$, that
$\kn{q}{r} > \eff - \dist{q}{r}$. By \ref{f:clock} and \ref{f:delivery} the
heartbeats of \rn{Tick} raise $\kn{q}{r}$ without bound unless $r$ is
itself permanently blocked. So define the \emph{demand relation}: a process
blocked at effective time $\eff$ on process $\proc{q}$ demands that each
$r \in \nbr{q}$ reach origin time strictly greater than
$\eff - \dist{q}{r}$.

Each step of this relation strictly decreases the demanded time, by at least
$\dist{q}{r} \ge 1$ (M2). A cyclic wait would therefore require a chain
$\eff > \eff_1 > \dots > \eff_k = \eff$, which is impossible. Since $\Ranks$
is finite, the demand chain from any blocked process is a finite acyclic chain
terminating at a process whose demand is vacuous, and that process makes progress
by \ref{f:clock}. Working back along the chain, every blocked process is
eventually released.
\end{proof}

\begin{remark}[Why (M2) is not negotiable]\label{rem:zero-cycle}
The proof consumes (M2) exactly once, essentially. Suppose $\dist{p}{q} = 0$
were permitted with its natural reading, ``visible in the same tick''. Then
$p$'s \rn{Publish} clause at $\eff$ reads $\kn{p}{q} > \eff$ and $q$'s reads
$\kn{q}{p} > \eff$: each must observe the other strictly beyond $\eff$ before
either may publish $\eff$, and neither can advance past $\eff$ until it has
read the other's version at $\eff$. The demand chain closes into a zero-weight
cycle and both block forever --- the smallest instance of the classical
zero-lookahead deadlock \cite{chandy1979} \devmark{dev:zero-distance}.
\end{remark}

\begin{theorem}[Barrier correctness]\label{thm:barrier-correct}
Under \cref{def:admissible} and \cref{res:collective,res:terminate}, the
three-round protocol terminates at every process, and when a process returns from
it, $\Pend{p} = \emptyset$ at every process and $\Net$ contains no
$\mathsf{PUT}$, $\mathsf{GETREQ}$ or $\mathsf{GETREP}$.
\end{theorem}

\begin{proof}[Proof sketch]
First round terminates because each process's outgoing set is finite and
\ref{f:delivery} completes the sends. When the first-round token returns to process $0$,
every process has passed \rn{Bar-One}, so no process will emit a further
$\mathsf{PUT}$ or $\mathsf{GETREQ}$. Second round sets $\mathit{drained}_p$, which
removes the neighbour clauses, so every remaining version publishes and
retires without further external input; every pending $\mathsf{iget}$ is then
served and every $\mathsf{oget}$ receives its reply, so $\Pend{p}$ drains.
When the second-round token returns to process $0$, this has happened everywhere. The exit
round propagates that conclusion.
\end{proof}

\begin{proposition}[Three rounds are necessary]\label{thm:barrier-three-rounds}
No two-round ring protocol establishes the postcondition of
\cref{thm:barrier-correct} at every process.
\end{proposition}

\begin{proof}
Two distinct global predicates must be established in sequence:
$\Phi_1$, ``no process will issue another operation'', and $\Phi_2$, ``no process
holds an incomplete request''. $\Phi_2$ cannot be evaluated locally before
$\Phi_1$ is known globally, because until every process has stopped issuing, a
process with an empty queue may still receive a request. A single ring traversal
establishes one predicate, and establishes it \emph{only at the initiator}:
when the token returns to process $0$, process $0$ knows $\Phi_1$, but process $\proc{i}$ with $i > 0$
knows only that processes $0..i$ had passed when the token crossed it.

Hence traversal one establishes $\Phi_1$ at process $0$, and traversal two both
propagates $\Phi_1$ and collects $\Phi_2$, again terminating the knowledge at
process $0$. A process $\proc{i}$ with $i > 0$ that exited after forwarding the second token would
do so without knowing $\Phi_2$ holds at processes $\proc{j}$ with $j > i$; such a process may still
be driving progress, and tearing down \gls{region} $i$ at that moment orphans any
reply still owed to it. A third traversal is required to announce $\Phi_2$.
\end{proof}

% ===========================================================================
\subsection{Refinement}
\label{sec:l1-refinement}

\begin{corollary}[$\Lom$ implements $\Lzm$]\label{cor:refinement}
For every deterministic program whose executions are legal, every admissible
$\Lom$ execution satisfies \cref{def:correctness}.
\end{corollary}

\begin{proof}
Clause (ii) is \cref{thm:read-soundness}. Clause (i) is
\cref{thm:deadlock-freedom} for the gets and \cref{thm:barrier-correct} for
\gls{termination}.
\end{proof}

The refinement map is the annotation on each rule in \cref{sec:l1-rules}.
Three aspects of $\Lim$ are abstracted away:

\begin{description}[leftmargin=1.1cm,style=nextline,font=\normalfont\bfseries]
\item[The allocation policy.] $\Lom$ fixes only the contract of
  \cref{def:alloc-ops} --- fresh handle, contiguous extent, disjoint from
  every live one, stable until freed --- and \cref{ts:alloc} leaves the choice
  of interval nondeterministic. Any policy resolving that nondeterminism
  refines the specification, so the scheme
  \code{dstm\_region\_malloc} happens to implement is not modelled. It
  determines which addresses a derived structure occupies, never what a
  \gls{cell} contains.

\item[MPI tags, buffers and packing.] The $\iota$ of \rn{Get-Remote} is
  the only piece of the messaging layer $\Lom$ retains, because reply matching
  is semantically visible. Packing, non-blocking send progress and the
  request status DFA of \srcrefr{request.h}{43}{46} refine the atomic
  \rn{Recv} and \rn{Serve} rules by \ref{f:delivery}.

\item[Session communicators.] \code{dstm\_start} duplicates a private
  communicator per \gls{session}, which makes $\Net$ per-session. $\Lom$ models one
  session, so this is exactly the assumption that $\Net$ starts empty.
\end{description}

\section{Synchronisation and Dependences}
\label{sec:sync}

The previous sections presented DSTM as a memory system. This one presents it
as a \emph{schedule}, in the vocabulary of dependence analysis and scheduling.
Nothing new is assumed: everything follows from
\crefrange{res:monotone}{res:terminate}. The payoff is that the results become
graph-theoretic and the design's structure becomes visible.

% ===========================================================================
\subsection{Bernstein conditions over versioned cells}
\label{sec:sync-bernstein}

Bernstein's conditions \cite{bernstein1966} state when two program fragments
may be reordered: when neither writes what the other touches. With footprints
as in \cref{def:footprint}:

\begin{definition}[Independence]\label{def:independence-ops}
Operations $o_1, o_2$ are \emph{independent}, written
$o_1 \disjoint o_2$, when
\[
  \underbrace{\fpw{o_1} \cap \fpr{o_2} = \emptyset}_{\text{no flow (RAW)}}
  \quad\wedge\quad
  \underbrace{\fpr{o_1} \cap \fpw{o_2} = \emptyset}_{\text{no anti (WAR)}}
  \quad\wedge\quad
  \underbrace{\fpw{o_1} \cap \fpw{o_2} = \emptyset}_{\text{no output (WAW)}} .
\]
\end{definition}

The decisive structural fact is that a footprint carries an effective time.

\begin{theorem}[Versioning is renaming]\label{thm:renaming}
If $\opeff{o_1} \ne \opeff{o_2}$ then $o_1 \disjoint o_2$, whatever their
\glspl{rank}, addresses and lengths.
\end{theorem}

\begin{proof}
By \cref{def:footprint}, $\fp{o_i} \subseteq \{\optgt{o_i}\} \times
\{\opeff{o_i}\} \times \Addr$. \Glspl{cell} carry their effective time as a
coordinate, so footprints at distinct \glspl{effective-time} are disjoint, and all
three intersections of \cref{def:independence-ops} are empty.
\end{proof}

\Cref{thm:renaming} is short but is why DSTM needs neither locks nor rollback.
In a conventional \gls{darray}, overwriting a location at tick $t+1$ creates
an output dependence on the write at tick $t$ and an anti-dependence on every
in-flight read of tick $t$; enforcing those is what a per-tick \gls{barrier}
does. DSTM removes them by \emph{renaming}: the tick index is part of the
location's name, so the two writes are not to the same location. This is the
manoeuvre of static single assignment \cite{cytron1991ssa} and array expansion
\cite{feautrier1991}, applied along a simulation-time axis.

What survives is only what happens \emph{within} one effective time.

\begin{theorem}[Residual dependences]\label{thm:residual-dep}
Let $\Hist$ be legal. Then for all $o_1 \ne o_2$ in $\Hist$, either
$o_1 \disjoint o_2$, or $o_1$ and $o_2$ are a put and a get with
$\optgt{o_1} = \optgt{o_2}$, $\opeff{o_1} = \opeff{o_2}$ and overlapping
addresses --- a flow dependence from the put to the get.
\end{theorem}

\begin{proof}
By \cref{thm:renaming} assume $\opeff{o_1} = \opeff{o_2} = \eff$ and
$\optgt{o_1} = \optgt{o_2}$, otherwise independence is immediate. Three cases.

\emph{Both \glspl{put}.} Then $\fpw{o_1} \cap \fpw{o_2} = \emptyset$ by
\cref{res:write-disjoint}, so there is no output dependence.

\emph{Both \glspl{get}.} $\fpw{o_i} = \emptyset$ for both, so all three intersections
are empty.

\emph{One put $w$, one get $g$.} The only possibly non-empty intersection is
$\fpw{w} \cap \fpr{g}$, a flow dependence. The anti-dependence
$\fpr{w} \cap \fpw{g}$ is empty because $\fpr{w} = \fpw{g} = \emptyset$ by
definition.
\end{proof}

\begin{corollary}[Normal form]\label{cor:dsa}
Under legality, a DSTM execution is in dynamic single assignment form: every
cell is written at most once, and the only dependences are flow dependences
from that unique write to its readers. Anti- and output dependences do not
occur at any granularity.
\end{corollary}

Two things eliminate them. Across effective times \cref{thm:renaming} does so
\emph{structurally} --- no restriction is involved and no program can
reintroduce them. Within one effective time \cref{res:write-disjoint}
eliminates output dependences and \cref{res:phase} anti-dependences, both
program restrictions a bad program can violate. DSTM's determinism is thus
partly free and partly earned.

% ===========================================================================
\subsection{The dependence graph}
\label{sec:sync-dag}

\begin{definition}[Dependence graph]\label{def:depgraph}
The \emph{dependence graph} $\depgraph(\Hist)$ has
\begin{itemize}[leftmargin=1.5em,itemsep=1pt]
  \item a node $\task{p}{t}$ for each task (\cref{def:task}) of $\Hist$, with
    node weight $\work{p}{t}$;
  \item a \emph{sequential} edge $\task{p}{t} \to \task{p}{t'}$ whenever $t <
    t'$ are consecutive \glspl{origin-time} used by $p$;
  \item a \emph{flow} edge $\task{p}{t} \to \task{r}{t'}$ whenever a put
    issued by $p$ at origin $t$ and a get issued by $r$ at origin $t'$ stand
    in the flow dependence of \cref{thm:residual-dep}.
\end{itemize}
\end{definition}

\begin{lemma}[Edge span]\label{lem:edge-span}
Let $\task{p}{t} \to \task{r}{t'}$ be a flow edge arising from a put with
\gls{delta} $\delta$ and a get with delta $\delta'$, both at effective time $\eff$
on owner $q$. Then
\[
  t' - t \;=\; \delta + |\delta'| ,
\]
and $t' - t \ge 1$ whenever $p \ne r$.
\end{lemma}

\begin{proof}
$t = \eff - \delta$ and $t' = \eff - \delta' = \eff + |\delta'|$, giving the
identity. For the bound, suppose $p \ne r$. If $p \ne q$ then
\cref{res:lookahead} and (M2) give $\delta \ge \dist{p}{q} \ge 1$. If $p = q$
then $r \ne q$, so the get is remote and the rule \textsc{Get-Remote} requires
$\delta' \le -1$, giving $|\delta'| \ge 1$. Either way $t' - t \ge 1$.
\end{proof}

\begin{theorem}[Schedulability]\label{thm:acyclic}
For a legal history, $\depgraph(\Hist)$ is acyclic.
\end{theorem}

\begin{proof}
Order nodes by $\task{p}{t} \prec \task{r}{t'}$ iff $t < t'$, or $t = t'$ and
the operations of $\task{p}{t}$ precede those of $\task{r}{t'}$ in the program
order of the common process. Sequential edges increase $t$ by construction. Flow
edges with $p \ne r$ increase $t$ strictly by \cref{lem:edge-span}. Flow edges
with $p = r$ have $t' - t = \delta + |\delta'| \ge 0$, and when $t' = t$ both
$\delta$ and $\delta'$ are zero, so the put and the get are local, at the same
effective time, on the same process --- and \cref{res:phase} orders the put
before the get in \gls{program-order}. Every edge therefore strictly increases
$\prec$, which is a strict partial order, so no cycle exists.
\end{proof}

\Cref{thm:acyclic} is the graph-theoretic form of \cref{thm:determinism}: the
induction used there is a topological order of $\depgraph$. It is also that of
\cref{thm:deadlock-freedom}, since a deadlock is a cycle in the wait-for
graph, itself a sub-graph of $\depgraph$ reversed.

\begin{remark}[Recovering the Chandy--Misra condition]\label{rem:chandy-misra}
\Cref{lem:edge-span} is the only place positivity is used, to show every
cross-rank edge spans at least one tick. Conservative PDES states the same
requirement as ``no cycle of zero lookahead'' \cite{chandy1979}; here it is a
property of $\cgraph$, holding by construction rather than by analysis, which
is what lets DSTM dispense with null messages.
\end{remark}

% ===========================================================================
\subsection{Push and pull are the same schedule}
\label{sec:sync-duality}

\begin{figure}[htbp]
\centering
\resizebox{\textwidth}{!}{%
\begin{tikzpicture}[
  x=20mm, y=-13mm,
  every node/.style={font=\scriptsize},
  tk/.style={draw,rounded corners=1.5pt,minimum width=13mm,minimum height=6.5mm,
             fill=resbg,draw=dstmblue,inner sep=1pt},
  vv/.style={draw,rounded corners=1.5pt,minimum width=13mm,minimum height=6.5mm,
             fill=white,draw=dstmgrey,dashed,inner sep=1pt},
  >={Stealth[round]},
]
  % ---- pull ----
  \node[font=\small\bfseries,text=dstmblue,anchor=west] at (-0.9,-0.9) {Pull};
  \node[anchor=east,text=dstmgrey] at (-0.35,0) {$p$};
  \node[anchor=east,text=dstmgrey] at (-0.35,1) {$r$};
  \node[tk] (pa) at (0,0) {$t$};
  \node[vv] (va) at (1,0) {$\ver{p}{t}$};
  \node[tk] (ra) at (2,1) {$t{+}D$};
  \draw[->,thick,dstmblue] (pa) -- node[above,inner sep=1.5pt]{$\opput\ \delta{=}0$} (va);
  \draw[->,thick,dashed,dstmred] (va) -- node[above right,inner sep=1.5pt]{$\opget\ \delta{=}{-}D$} (ra);

  % ---- push ----
  \node[font=\small\bfseries,text=dstmblue,anchor=west] at (-0.9,2.1) {Push};
  \node[anchor=east,text=dstmgrey] at (-0.35,3) {$p$};
  \node[anchor=east,text=dstmgrey] at (-0.35,4) {$r$};
  \node[tk] (pb) at (0,3) {$t$};
  \node[vv] (vb) at (1,4) {$\ver{r}{t{+}D}$};
  \node[tk] (rb) at (2,4) {$t{+}D$};
  \draw[->,thick,dstmblue] (pb) -- node[above right,inner sep=1.5pt]{$\opput\ \delta{=}{+}D$} (vb);
  \draw[->,thick,dashed,dstmred] (vb) -- node[above,inner sep=1.5pt]{$\opget\ \delta{=}0$} (rb);

  % ---- common conclusion ----
  \draw[decorate,decoration={brace,amplitude=4pt,mirror},dstmamber,thick]
       (2.42,-0.35) -- (2.42,4.35)
       node[midway,right=6pt,text=dstmamber,align=left,font=\scriptsize]
       {consumer tick is\\$t + D$ either way};
\end{tikzpicture}}
\caption{The two idioms of \cref{sec:exec-idioms}. Solid edges are writes,
dashed edges reads, dashed boxes versions. The version they route through
differs --- $p$'s own region versus $r$'s --- but the tick at which the value
reaches the consumer is $t + \dist{p}{r}$ in both cases.}
\label{fig:push-pull}
\end{figure}

\begin{theorem}[Pipeline depth]\label{thm:pipeline-depth}
Let $p$ produce a value at origin time $t$ and let $r$ consume it, with
$q \in \{p, r\}$ the owner of the mediating version, $\dist{p}{r} = D$.
\begin{enumerate}[label=(\roman*),leftmargin=1.4cm,itemsep=1pt]
  \item In the \glslink{pull-model}{pull idiom} ($q = p$, $\delta = 0$, $\delta' = -D$) the consumer
    tick is $t' = t + D$.
  \item In the \glslink{push-model}{push idiom} ($q = r$, $\delta = +D$, $\delta' = 0$) the consumer
    tick is $t' = t + D$.
\end{enumerate}
Moreover $t' = t + D$ is optimal: no legal choice of $\delta, \delta'$ gives
$t' < t + D$.
\end{theorem}

\begin{proof}
Both cases are \cref{lem:edge-span}: $t' - t = \delta + |\delta'|$, which is
$0 + D$ and $D + 0$ respectively. For optimality, if $q = p$ then
$\delta \ge 0$ and, since $r \ne q$, \cref{res:staleness} bounds
$|\delta'| \le \dist{r}{p} = D$ while \textsc{Get-Remote} forces
$|\delta'| \ge 1$; the minimum of $\delta + |\delta'|$ subject to the reader
actually reaching $p$'s version is attained at $\delta = 0$, and the consumer
tick can be no earlier than the tick at which the reader may legally look back
to $t$, namely $t + |\delta'|$ with $|\delta'| \le D$. Symmetrically if
$q = r$ then $\delta \ge \dist{p}{r} = D$ and $\delta' = 0$, so
$t' - t \ge D$. In both cases $t' - t \ge \min(D, D) = D$ once the constraint
that the value must actually cross from $p$ to $r$ is imposed.
\end{proof}

\begin{remark}[Push or pull is a cost decision]
The corollary is easy to miss: the choice between push and pull is never
semantic and never affects how far processes may drift. It is a
communication-cost decision. Push moves one message per producer--consumer
pair per tick, pull two; push is the better default and is what both benchmark
applications in \cite{dstm2026} use. Pull earns its cost when the consumer
reads only occasionally, moving data on demand.
\end{remark}

% ===========================================================================
\subsection{Slack, drift and the critical path}
\label{sec:sync-slack}

Since $\depgraph$ is a weighted DAG, the standard scheduling quantities are
available.

\begin{definition}[Earliest start, slack]\label{def:slack}
For a node $u$ of $\depgraph$, let $\mathit{est}(u)$ be the length of the
longest weighted path ending at $u$ (excluding $u$'s own weight), and
$\mathit{lst}(u)$ the latest start compatible with the overall makespan. The
\emph{slack} is $\slack{}{u} \defeq \mathit{lst}(u) - \mathit{est}(u)$.
\end{definition}

Drift (\cref{def:drift}) realises slack: a process runs ahead exactly as far
as its tasks have slack, and stops --- in a blocking get --- on reaching a
task with none. A barrier-synchronised execution zeroes all slack at every
tick boundary, which is why it is slower and why its cost is set by the
slowest process.

\begin{theorem}[Makespan]\label{thm:makespan}
Let $\cpath(\Hist)$ be the maximum, over directed paths $\pi$ in
$\depgraph(\Hist)$, of $\sum_{u \in \pi} \work{}{u}$. Then, under
\cref{def:admissible} and with zero communication cost,
\[
  \makespan_{\mathrm{DSTM}} \;=\; \cpath(\Hist)
  \qquad\text{and}\qquad
  \makespan_{\mathrm{BSP}} \;=\; \sum_{t} \max_{p} \work{p}{t},
\]
and $\makespan_{\mathrm{DSTM}} \le \makespan_{\mathrm{BSP}}$.
\end{theorem}

\begin{proof}
The first equality is the standard characterisation of a greedy schedule on a
DAG with unbounded processors, available because \cref{thm:acyclic} gives
acyclicity and each process executes its own tasks in order. For the inequality,
observe from the proof of \cref{thm:acyclic} that every edge strictly
increases the origin time or, at equal origin time, the program-order
position; hence a directed path visits at most one node per origin time $t$.
Therefore
\[
  \sum_{u \in \pi} \work{}{u}
  \;\le\; \sum_{t} \max_{p} \work{p}{t}
  \;=\; \makespan_{\mathrm{BSP}} . \qedhere
\]
\end{proof}

\begin{remark}[Where the measured speedup comes from]
The gap between the two expressions in \cref{thm:makespan} is what the
evaluation in \cite{dstm2026} reports. It is large when per-tick load is
imbalanced but not persistently so: when the expensive process differs from
tick to tick, $\sum_t \max_p \work{p}{t}$ sums a sequence of maxima while the
critical path weaves between processes. It vanishes when one process is
uniformly slowest, the critical path then running down that row.
\end{remark}

% ===========================================================================
\subsection{Memory as a function of drift}
\label{sec:sync-space}

Versioning trades memory for synchronisation. The trade is bounded, which is
why the scheme is practical.

\begin{theorem}[Space bound]\label{thm:space-bound}
In any reachable configuration, the number of live versions at process $\proc{p}$
satisfies
\[
  \bigl|\{\eff : \Vst{p}{\eff} \ne \Absent\}\bigr|
  \;\le\; H_p - L_p + 1,
\]
where $H_p$ is the largest effective time targeted by any put issued so far
towards $p$, and
\[
  L_p \;\defeq\; \min_{q \in \nbr{p}} \bigl(\kn{p}{q} - \dist{p}{q}\bigr),
  \qquad L_p \defeq \clk{p} \text{ when } \nbr{p} = \emptyset .
\]
\end{theorem}

\begin{proof}
A version exists only if created by \textsc{Apply}, so no live version exceeds
$H_p$. For the lower end, any version with $\eff < L_p$ satisfies
$\kn{p}{q} - \dist{p}{q} > \eff$ for every $q \in \nbr{p}$; together with
$\clk{p} \ge \eff$ (which holds since $L_p \le \clk{p}$ by
\cref{inv:l1-clock}) and the absence of pending gets --- guaranteed by
\cref{thm:discard-safety}, since such a get would be illegal --- the premises
of \textsc{Retire} hold, and \textsc{Evict} follows by \ref{f:progress}.
\end{proof}

\begin{corollary}[Memory is linear in drift]\label{cor:space-drift}
Suppose the execution has bounded drift $\Theta$ (\cref{def:bounded-drift}),
the horizon bound of \cref{res:horizon} is $\Theta_{\max}$, and heartbeats are
delivered promptly enough that $\kn{p}{q} \ge \clk{q} - \lambda$ for a
constant $\lambda$. Then
\[
  \bigl|\{\eff : \Vst{p}{\eff} \ne \Absent\}\bigr|
  \;\le\; 2\Theta + \Theta_{\max} + \diam_p + \lambda + 1 .
\]
In particular the live-version count is independent of the length of the run.
\end{corollary}

\begin{proof}
$H_p \le \max_{q \in \nbr{p} \cup \{p\}} \clk{q} + \Theta_{\max}
 \le \clk{p} + \Theta + \Theta_{\max}$ by bounded drift and
\cref{res:horizon}. And
$L_p \ge \min_q (\clk{q} - \lambda) - \diam_p
 \ge \clk{p} - \Theta - \lambda - \diam_p$. Subtract.
\end{proof}

\Cref{cor:space-drift} is the formal content of the claim in \cite{dstm2026}
that discarding \glspl{version} ``is a memory management optimization $\dots$
without changing the semantics of the program''. Both halves are theorems:
semantics-preservation is \cref{thm:discard-safety}, the bound is
\cref{cor:space-drift}. Each term says something operational --- $2\Theta$
that memory is the price of \gls{drift}, $\Theta_{\max}$ that writing far into
the future costs, $\diam_p$ that distant neighbours cost, $\lambda$ that a
lazy \gls{heartbeat} implementation costs memory rather than correctness.

% ===========================================================================
\subsection{Summary}
\label{sec:sync-summary}

\begin{table}[htbp]
\centering
\small
\begin{tabularx}{\textwidth}{@{}>{\raggedright\arraybackslash}p{4.1cm}
                                >{\raggedright\arraybackslash}X l@{}}
\toprule
\textbf{DSTM mechanism} & \textbf{Dependence-theoretic reading}
 & \textbf{Result} \\
\midrule
Effective time in the cell identity
 & renaming; kills WAR and WAW across ticks
 & \cref{thm:renaming} \\
\Cref{res:write-disjoint}
 & dynamic single assignment within a tick
 & \cref{cor:dsa} \\
\Cref{res:phase}
 & write-phase before read-phase per version; kills WAR within a tick
 & \cref{thm:residual-dep} \\
Delta of an operation
 & weight of a producer--consumer edge
 & \cref{lem:edge-span} \\
(M2), $\dist{p}{q} \ge 1$
 & every cycle of $\cgraph$ has positive weight
 & \cref{thm:acyclic} \\
Distance matrix
 & declared lookahead; pipeline depth
 & \cref{thm:pipeline-depth} \\
Process drift
 & slack in the schedule
 & \cref{def:slack} \\
Absence of a per-tick barrier
 & critical path instead of sum-of-maxima
 & \cref{thm:makespan} \\
Version retirement
 & liveness analysis; storage reclamation
 & \cref{thm:space-bound} \\
Three-round ring barrier
 & termination detection
 & \cref{thm:barrier-correct} \\
\bottomrule
\end{tabularx}
\caption{The dictionary between DSTM's mechanisms and their
dependence-theoretic content.}
\label{tab:sync-dictionary}
\end{table}

